\Article{61}
So bequem diese Methode bei vorhandenen Tafeln ist, dürfen wir doch nicht vergessen, dass sie indirekt ist. Es lohnt sich also zu untersuchen, was direkte Methoden leisten können. Hier geben wir die Ergebnisse, die sich aus dem Vorhergehenden gewinnen lassen; alles Weitere, das tiefere Überlegungen verlangt, bleibt Abschnitt VIII vorbehalten. Wir beginnen mit dem einfachsten Fall $A=1$, suchen also die Wurzeln der Kongruenz $x^n\equiv1\,(\operatorname{mod.}p)$. Nimmt man eine beliebige primitive Wurzel als Basis, so muss $n\operatorname{Ind.}x\equiv0\,(\operatorname{mod.}p-1)$ sein. Ist $n$ zu $p-1$ teilerfremd, so hat diese Kongruenz nur eine Wurzel, nämlich $\operatorname{Ind.}x\equiv0\,(\operatorname{mod.}p-1)$. Also hat \emph{in diesem Fall} ${}^{n}\!\surd1\,(\operatorname{mod.}p)$ einen einzigen Wert, nämlich $\equiv1$. Haben dagegen die Zahlen $n,p-1$ den größten gemeinsamen Teiler $\delta$, so lautet die vollständige Lösung der Kongruenz $n\operatorname{Ind.}x\equiv0\,(\operatorname{mod.}p-1)$: $\operatorname{Ind.}x\equiv0\,(\operatorname{mod.}\frac{p-1}{\delta})$ (siehe Artikel 29). \emph{Das heißt,} $\operatorname{Ind.}x$ muss nach dem Modul $p-1$ einer der Zahlen
\[0,\frac{p-1}{\delta}.\ \frac{2(p-1)}{\delta},\frac{3(p-1)}{\delta},\GaussFourDots\frac{(\delta-1)(p-1)}{\delta}\]
kongruent sein oder $\delta$ nach dem Modul $p-1$ inkongruente Werte haben. Somit hat auch $x$ in diesem Fall $\delta$ verschiedene Werte (inkongruent nach dem Modul
\SourceLeaf{55}$p$). Daraus ergibt sich, dass auch der Ausdruck ${}^{\delta}\!\surd {1}$ genau $\delta$ verschiedene Werte hat, deren Indizes mit den eben angegebenen vollständig übereinstimmen. Der Ausdruck ${}^{\delta}\!\surd {1}\,(\operatorname{mod.}p)$ ist also mit ${}^{n}\!\surd {1}\,(\operatorname{mod.}p)$ gleichwertig: Die Kongruenz $x^\delta\equiv1\,(\operatorname{mod.}p)$ hat dieselben Wurzeln wie $x^n\equiv1\,(\operatorname{mod.}p)$. Die erste hat aber einen niedrigeren Grad, falls $\delta$ und $n$ ungleich sind.

\emph{Beispiel.} ${}^{15}\!\surd {1}\,(\operatorname{mod.}19)$ hat drei Werte, da $3$ der größte gemeinsame Teiler der Zahlen $15,18$ ist. Diese sind zugleich die Werte von ${}^{3}\!\surd {1}\,(\operatorname{mod.}19)$, nämlich $1,7,11$.

\Article{62}
Diese Reduktion hat also den Vorteil, dass wir nur noch diejenigen Kongruenzen der Form $x^n\equiv1$ lösen müssen, bei denen $n$ ein Teiler von $p-1$ ist. Weiter unten zeigen wir, dass sich Kongruenzen dieser Form stets noch weiter im Grad herabsetzen lassen; die bisherigen Ergebnisse reichen dazu allerdings nicht aus. Einen Fall können wir schon hier erledigen, nämlich $n=2$. Die Werte des Ausdrucks ${}^{2}\!\surd {1}$ sind offenbar $+1$ und $-1$, denn mehr als zwei kann er nicht haben. Dabei sind $+1$ und $-1$ stets inkongruent, außer wenn der Modul $=2$ ist. In diesem Fall kann ${}^{2}\!\surd {1}$ offensichtlich nur einen Wert haben. Daraus folgt, dass $+1$ und $-1$ auch die Werte von ${}^{2m}\!\surd {1}$ sind, wenn $m$ zu $\frac{p-1}{2}$ teilerfremd ist. Dies gilt immer dann, wenn der Modul so beschaffen ist, dass $\frac{p-1}{2}$ eine Primzahl ist (außer wenn $p-1=2m$ gilt; dann sind sämtliche Zahlen $1,2,3\GaussFourDots p-1$ Wurzeln), \emph{zum Beispiel} für $p=3,5,7,11,23,47,59,83,107$ usw. Als Folgerung sei hier festgehalten, dass der Index von $-1$ stets $\equiv\frac{p-1}{2}\,(\operatorname{mod.}p-1)$ ist, ganz gleich, welche primitive Wurzel als Basis gewählt wird. Denn $2\operatorname{Ind.}(-1)\equiv0\,(\operatorname{mod.}p-1)$. Also ist $\operatorname{Ind.}(-1)$ entweder $\equiv0$ oder $\equiv\frac{p-1}{2}\,(\operatorname{mod.}p-1)$. Nun ist aber $0$ immer der Index von $+1$, und $+1$ und $-1$ müssen stets verschiedene Indizes haben (abgesehen vom Fall $p=2$, den wir hier nicht weiter beachten müssen).

\Article{63}
In Artikel 60 haben wir gezeigt, dass der Ausdruck ${}^{n}\!\surd {A}\,(\operatorname{mod.}p)$ genau $\delta$ verschiedene Werte oder gar keinen hat, wenn $\delta$ der größte gemeinsame Teiler der Zahlen $n,p-1$ ist. Wie wir eben die Gleichwertigkeit von ${}^{n}\!\surd {A}$ und ${}^{\delta}\!\surd {A}$ für $A\equiv1$ gezeigt haben, beweisen wir nun allgemeiner, dass sich der Ausdruck ${}^{n}\!\surd {A}$ stets auf einen anderen Ausdruck ${}^{\delta}\!\surd {B}$ zurückführen lässt, mit dem er gleichwertig ist. Bezeichnen wir einen beliebigen Wert des ersten mit $x$, so ist $x^n\equiv A$. Nun
\SourceLeaf{56}sei $t$ irgendein Wert von $\frac{\delta}{n}\,(\operatorname{mod.}p-1)$; dass dieser Ausdruck wirkliche Werte hat, ergibt sich aus Artikel 31. Dann ist $x^{tn}\equiv A^t$, aber $x^{tn}\equiv x^\delta$ wegen $tn\equiv\delta\,(\operatorname{mod.}p-1)$. Also ist $x^\delta\equiv A^t$, und jeder Wert von ${}^{n}\!\surd {A}$ ist somit auch ein Wert von ${}^{\delta}\!\surd {A^t}$. Sobald ${}^{n}\!\surd {A}$ wirkliche Werte hat, ist dieser Ausdruck daher mit ${}^{\delta}\!\surd {A^t}$ vollständig gleichwertig, da er weder andere noch weniger Werte hat. Hat dagegen ${}^{n}\!\surd {A}$ keinen wirklichen Wert, so kann ${}^{\delta}\!\surd {A^t}$ trotzdem wirkliche Werte haben.

\emph{Beispiel.} Sucht man die Werte des Ausdrucks ${}^{21}\!\surd {2}\,(\operatorname{mod.}31)$, so ist der größte gemeinsame Teiler von $21$ und $30$ gleich $3$, und ein Wert von $\frac{3}{21}\,(\operatorname{mod.}30)$ ist $3$. Hat also ${}^{21}\!\surd {2}$ wirkliche Werte, so ist dieser Ausdruck mit ${}^{3}\!\surd {2^3}$ oder ${}^{3}\!\surd {8}$ gleichwertig. Tatsächlich findet man, dass die Werte des letzteren Ausdrucks, nämlich $2,10.19$, auch dem ersten genügen.

\Article{64}
Damit wir diese Rechnung nicht womöglich vergeblich ausführen, brauchen wir eine Regel, die sofort entscheidet, ob ${}^{n}\!\surd {A}$ wirkliche Werte hat oder nicht. Mit einer Indextafel ist dies leicht: Nach Artikel 60 gibt es wirkliche Werte, wenn der Index von $A$ bei einer beliebigen primitiven Wurzel als Basis durch $\delta$ teilbar ist; andernfalls gibt es keine. Man kann das aber auch ohne eine solche Tafel feststellen. Setzen wir nämlich den Index von $A$ gleich $k$, so ist, falls dieser durch $\delta$ teilbar ist, auch $\frac{k(p-1)}{\delta}$ durch $p-1$ teilbar, und umgekehrt. Nun ist aber der Index der Zahl $A^{\frac{p-1}{\delta}}$ gleich $\frac{k(p-1)}{\delta}$. Hat also ${}^{n}\!\surd {A}\,(\operatorname{mod.}p)$ wirkliche Werte, so ist $A^{\frac{p-1}{\delta}}$ der Eins kongruent, andernfalls nicht. Im Beispiel des vorigen Artikels gilt etwa $2^{10}=1024\equiv1\,(\operatorname{mod.}31)$, woraus folgt, dass ${}^{21}\!\surd {2}\,(\operatorname{mod.}31)$ wirkliche Werte hat.
Ebenso erfahren wir daraus, dass ${}^{2}\!\surd {-1}\,(\operatorname{mod.}p)$ immer zwei wirkliche Werte hat, wenn $p$ von der Form $4m+1$ ist, dagegen keinen, wenn $p$ von der Form $4m+3$ ist; denn $(-1)^{2m}=1$ und $(-1)^{2m+1}=-1$. Diesen schönen Satz, den man gewöhnlich so ausspricht: \emph{Ist $p$ eine Primzahl der Form $4m+1$, so lässt sich eine Quadratzahl $aa$ finden, sodass $aa+1$ durch $p$ teilbar ist; ist dagegen $p$ von der Form $4m-1$, so gibt es eine solche Quadratzahl nicht}, hat Euler auf diese Weise bewiesen, \emph{Comm. nov. Acad. Petrop. T.} XVIII, \emph{p.} 112, für das Jahr 1773. Einen anderen Beweis hatte er schon viel früher gegeben, \emph{Comm. nov. T.} V, \emph{p.} 5, erschienen 1760. In einer früheren Abhandlung, \emph{Comm. nov. T.} IV, \emph{p.} 25, hatte er die Sache noch nicht
\SourceLeaf{57}abgeschlossen. Später gab auch La Grange einen Beweis des Satzes, \emph{Nouveaux Mém. de l'Ac. de Berlin A.} 1775, \emph{p.} 342. Einen weiteren Beweis geben wir im nächsten Abschnitt, wo dieser Gegenstand eigens behandelt wird.

\Article{65}
Wir haben gezeigt, wie sich alle Ausdrücke ${}^{n}\!\surd {A}\,(\operatorname{mod.}p)$ auf solche zurückführen lassen, bei denen $n$ ein Teiler von $p-1$ ist, und haben ein Kriterium dafür gefunden, ob wirkliche Werte existieren. Nun betrachten wir solche Ausdrücke ${}^{n}\!\surd {A}\,(\operatorname{mod.}p)$, bei denen $n$ ein Teiler von $p-1$ ist, genauer. Zuerst zeigen wir, wie die einzelnen Werte eines Ausdrucks miteinander zusammenhängen; dann geben wir einige Verfahren an, mit denen man oft einen Wert finden kann.

\emph{Erstens:} Ist $A\equiv1$ und $r$ einer der $n$ Werte von ${}^{n}\!\surd {1}\,(\operatorname{mod.}p)$, gilt also $r^n\equiv1\,(\operatorname{mod.}p)$, so sind auch alle Potenzen von $r$ Werte dieses Ausdrucks. Dabei gibt es so viele verschiedene Werte, wie der Exponent angibt, zu dem $r$ gehört (Artikel 48). Gehört also der Wert $r$ zum Exponenten $n$, so umfassen die Potenzen von $r$, nämlich $r,r^2,r^3\GaussFourDots r^n$ (wobei man die letzte durch die \emph{Eins} ersetzen kann), alle Werte des Ausdrucks ${}^{n}\!\surd {1}\,(\operatorname{mod.}p)$. Welche Hilfsmittel es gibt, um solche zum Exponenten $n$ gehörigen Werte zu finden, erläutern wir in Abschnitt VIII ausführlicher.

\emph{Zweitens.} Ist $A$ der Eins inkongruent und ein Wert von ${}^{n}\!\surd {A}\,(\operatorname{mod.}p)$, etwa $z$, bekannt, so erhält man daraus die übrigen wie folgt. Die Werte von ${}^{n}\!\surd {1}$ seien

\[1,\ r,\ r^2\ldots r^{n-1}\]

wie eben gezeigt. Dann sind sämtliche Werte von ${}^{n}\!\surd {A}$ die folgenden:

\[z,\ zr,\ zr^2\GaussFourDots zr^{n-1}\]

Sie alle genügen der Kongruenz $x^n\equiv A$. Setzt man nämlich einen beliebigen von ihnen $\equiv zr^k$, so ist seine Potenz mit dem Exponenten $n$, nämlich $z^n r^{nk}$, wegen $r^n\equiv1$ und $z^n\equiv A$ der Zahl $A$ kongruent. Dass sie alle verschieden sind, ergibt sich leicht aus Artikel 23. Mehr Werte als diese, deren Anzahl $n$ ist, kann der Ausdruck ${}^{n}\!\surd {A}$ aber nicht haben. Ist also \emph{zum Beispiel} ein Wert von ${}^{2}\!\surd {A}$ gleich $z$, so ist der andere $-z$. Schließlich
\SourceLeaf{58}müssen wir daraus folgern, dass man nicht alle Werte von ${}^{n}\!\surd {A}$ finden kann, ohne zugleich sämtliche Werte von ${}^{n}\!\surd {1}$ zu kennen.

\Article{66}
Als Zweites wollten wir zeigen, in welchem Fall sich ein Wert von ${}^{n}\!\surd {A}\,(\operatorname{mod.}p)$ direkt finden lässt; dabei setzen wir voraus, dass $n$ ein Teiler von $p-1$ ist. Dies ist möglich, wenn ein Wert einer Potenz von $A$ kongruent ist. Da dieser Fall recht oft vorkommt, lohnt es sich, etwas dabei zu verweilen. Ein solcher Wert, \emph{falls es einen gibt,} sei $z$; es gelte also $z\equiv A^k$ und $A\equiv z^n\,(\operatorname{mod.}p)$. Daraus folgt $A\equiv A^{kn}$. Kennt man also eine Zahl $k$, für die $A\equiv A^{kn}$ gilt, so ist $A^k$ ein gesuchter Wert. Diese Bedingung ist aber gleichbedeutend mit $1\equiv kn\,(\operatorname{mod.}t)$, wobei $t$ den Exponenten bezeichnet, zu dem $A$ gehört (Artikel 46, 48). Damit diese Kongruenz lösbar ist, muss $n$ zu $t$ teilerfremd sein. In diesem Fall gilt $k\equiv\frac{1}{n}\,(\operatorname{mod.}t)$. Haben dagegen $t$ und $n$ einen gemeinsamen Teiler außer der Eins, so kann kein Wert $z$ einer Potenz von $A$ kongruent sein.

\Article{67}
Für diese Lösung müssen wir allerdings $t$ kennen. Sehen wir also, wie man vorgehen kann, wenn diese Zahl unbekannt ist. Zunächst ist leicht einzusehen, dass $t$ ein Teiler von $\frac{p-1}{n}$ sein muss, sofern ${}^{n}\!\surd {A}\,(\operatorname{mod.}p)$ wirkliche Werte hat, wie wir hier stets voraussetzen. Sei nämlich $y$ ein beliebiger Wert. Dann gilt sowohl $y^{p-1}\equiv1$ als auch $y^n\equiv A\,(\operatorname{mod.}p)$. Erhebt man beide Seiten der letzteren Kongruenz zur Potenz mit dem Exponenten $\frac{p-1}{n}$, so erhält man $A^{\frac{p-1}{n}}\equiv1$. Damit ist $\frac{p-1}{n}$ durch $t$ teilbar (Artikel 48). Ist nun $\frac{p-1}{n}$ zu $n$ teilerfremd, so lässt sich die Kongruenz $kn\equiv1$ des vorigen Artikels auch nach dem Modul $\frac{p-1}{n}$ lösen. Ein Wert von $k$, der die Kongruenz nach diesem Modul erfüllt, erfüllt sie offenbar auch nach dem Modul $t$, der ein Teiler von $\frac{p-1}{n}$ ist (Artikel 5). Damit ist das Gesuchte gefunden. Ist dagegen $\frac{p-1}{n}$ zu $n$ nicht teilerfremd, so entferne man sämtliche Primfaktoren von $\frac{p-1}{n}$, die zugleich $n$ teilen, aus $\frac{p-1}{n}$. So erhalten wir eine Zahl $\frac{p-1}{nq}$, die zu $n$ teilerfremd ist. Dabei bezeichnet $q$ das Produkt aller entfernten Primfaktoren. Gilt nun die Bedingung aus dem vorigen Artikel, dass $t$ zu $n$ teilerfremd ist, so ist auch $t$ zu $q$ teilerfremd und teilt somit ebenfalls $\frac{p-1}{nq}$. Löst man also die Kongruenz $kn\equiv1\,(\operatorname{mod.}\frac{p-1}{nq})$ (was möglich ist, da $n$ zu $\frac{p-1}{nq}$ teilerfremd ist), so erfüllt der Wert von $k$ die Kongruenz auch nach dem Modul
\SourceLeaf{59}$t$, wie verlangt. Das ganze Verfahren besteht darin, eine Zahl zu finden, die an die Stelle des unbekannten $t$ treten kann. Wir müssen uns aber gut merken, dass wir, wenn $\frac{p-1}{n}$ zu $n$ nicht teilerfremd ist, die Bedingung des vorigen Artikels vorausgesetzt haben. Gilt sie nicht, so sind alle Schlüsse falsch. Finden wir beim unbedachten Befolgen der angegebenen Regeln für $z$ einen Wert, dessen Potenz mit dem Exponenten $n$ der Zahl $A$ nicht kongruent ist, so zeigt das, dass die Bedingung nicht gilt und dieses Verfahren überhaupt nicht anwendbar ist.

\Article{68}
Aber auch in diesem Fall kann die bereits geleistete Arbeit oft nützlich sein. Es lohnt sich zu untersuchen, wie dieser falsche Wert mit den richtigen zusammenhängt. Nehmen wir also an, die Zahlen $k,z$ seien regelgerecht bestimmt, aber $z^n$ sei nicht $\equiv A\,(\operatorname{mod.}p)$. Können wir nun die Werte von ${}^{n}\!\surd{\frac{A}{z^n}}\,(\operatorname{mod.}p)$ bestimmen, so erhalten wir durch Multiplikation jedes einzelnen mit $z$ die Werte von ${}^{n}\!\surd{A}$. Ist nämlich $v$ ein Wert von ${}^{n}\!\surd{\frac{A}{z^n}}$, so gilt $(vz)^n\equiv A$. Der Ausdruck ${}^{n}\!\surd{\frac{A}{z^n}}$ ist aber insofern einfacher als ${}^{n}\!\surd{A}$, als $\frac{A}{z^n}\,(\operatorname{mod.}p)$ meistens zu einem kleineren Exponenten gehört als $A$. Genauer: Ist der größte gemeinsame Teiler der Zahlen $t,q$ gleich $d$, so gehört $\frac{A}{z^n}\,(\operatorname{mod.}p)$ zum Exponenten $d$. Dies beweist man so. Setzt man den Wert von $z$ ein, so erhält man $\frac{A}{z^n}\equiv\frac{1}{A^{kn-1}}\,(\operatorname{mod.}p)$. Nun ist $kn-1$ durch $\frac{p-1}{nq}$ teilbar (voriger Artikel), $\frac{p-1}{n}$ aber durch $t$ (ebenda), oder $\frac{p-1}{nd}$ durch $\frac{t}{d}$. Nach Voraussetzung ist jedoch $\frac{t}{d}$ zu $\frac{q}{d}$ teilerfremd. Also ist auch $\frac{p-1}{nd}$ durch $\frac{tq}{dd}$ teilbar, oder $\frac{p-1}{nq}$ durch $\frac{t}{d}$, und damit auch $kn-1{}^{\prime}$ durch $\frac{t}{d}$ sowie $(kn-1)d$ durch $t$. Daher ist $A^{(kn-1)d}\equiv1\,(\operatorname{mod.}p)$. Daraus ergibt sich leicht, dass $\frac{A}{z^n}$, zur Potenz mit dem Exponenten $d$ erhoben, der Eins kongruent ist. Dass $\frac{A}{z^n}$ nicht zu einem kleineren Exponenten als $d$ gehören kann, lässt sich zwar leicht beweisen; da wir es aber für unseren Zweck nicht brauchen, verweilen wir nicht dabei. Wir können also sicher sein, dass $\frac{A}{z^n}\,(\operatorname{mod.}p)$ stets zu einem kleineren Exponenten gehört als $A$, mit einer einzigen Ausnahme: wenn $t$ ein Teiler von $q$ ist, also $d=t$.

Was bringt es uns aber, dass $\frac{A}{z^n}$ zu einem kleineren Exponenten gehört als $A$? Es gibt mehr Zahlen, die als $A$ auftreten können, als solche, die als $\frac{A}{z^n}$ auftreten können. Wenn mehrere solche Ausdrücke ${}^{n}\!\surd{A}$ nach demselben Modul auszuwerten sind, gewinnen wir dadurch, dass wir mehrere aus derselben Quelle erhalten können. So können wir \emph{zum Beispiel} stets mindestens einen Wert von ${}^{2}\!\surd{A}\,(\operatorname{mod.}29)$ bestimmen,
\SourceLeaf{60}wenn nur die Werte von ${}^{2}\!\surd{-1}\,(\operatorname{mod.}29)$ (nämlich $\pm12$) bekannt sind. Aus dem vorigen Artikel ergibt sich leicht, dass sich ein Wert solcher Ausdrücke stets direkt bestimmen lässt, wenn $t$ ungerade ist, und dass $d$ dann $=2$ ist, wenn $t$ gerade ist. Außer $-1$ gehört aber keine Zahl zum Exponenten $2$.

\emph{Beispiele.} Gesucht ist ${}^{3}\!\surd{31}\,(\operatorname{mod.}37)$. Hier ist $p-1=36$, $n=3$, $\frac{p-1}{3}=12$, also $q=3$. Es muss daher $3k\equiv1\,(\operatorname{mod.}4)$ gelten, was man mit $k=3$ erreicht. Damit ist $z\equiv31^3\,(\operatorname{mod.}37)\equiv6$, und tatsächlich findet man $6^3\equiv31\,(\operatorname{mod.}37)$. Sind die Werte von ${}^{3}\!\surd{1}\,(\operatorname{mod.}37)$ bekannt, so lassen sich auch die übrigen Werte von ${}^{3}\!\surd{6}$ bestimmen. Jene sind $1,10,26$. Multipliziert man $6$ mit ihnen, so erhält man die übrigen Werte $\equiv23$ und $8$.

Sucht man dagegen einen Wert von ${}^{2}\!\surd{3}\,(\operatorname{mod.}37)$, so ist $n=2$, $\frac{p-1}{n}=18$, also $q=2$. Daher muss $2k\equiv1\,(\operatorname{mod.}9)$ gelten, woraus $k\equiv5\,(\operatorname{mod.}9)$ folgt. Somit ist $z\equiv3^5\equiv21\,(\operatorname{mod.}37)$. Doch $21^2$ ist nicht $\equiv3$, sondern $\equiv34$. Nun ist aber $\frac{3}{34}\,(\operatorname{mod.}37)\equiv-1$ und ${}^{2}\!\surd{-1}\,(\operatorname{mod.}37)\equiv\pm6$, woraus man die richtigen Werte $\pm6\,.\,21\equiv\pm15$ erhält.

Das ist etwa alles, was wir hier über das Auswerten solcher Ausdrücke mitteilen können. Direkte Methoden werden offenbar oft recht langwierig. Dieser Nachteil betrifft aber fast alle direkten Methoden in der Zahlentheorie; deshalb wollten wir es nicht versäumen zu zeigen, was sie hier leisten können. Auch ist anzumerken, dass es nicht zu unserem Vorhaben gehört, die besonderen Kunstgriffe einzeln zu erläutern, die sich Geübten häufig von selbst anbieten.

\begin{center}\emph{Zusammenhang der Indizes in verschiedenen Systemen.}\end{center}

\Article{69}
Wir kehren nun zu den Wurzeln zurück, die wir primitive genannt haben. Wir haben gezeigt, dass bei jeder primitiven Wurzel als Basis alle Zahlen, deren Indizes zu $p-1$ teilerfremd sind, ebenfalls primitive Wurzeln sind, und keine anderen. Damit ergibt sich zugleich unmittelbar die Anzahl der primitiven Wurzeln; siehe Artikel 53. Welche primitive Wurzel wir als Basis wählen, steht uns im Allgemeinen frei. Also kann es auch hier, wie beim Rechnen mit Logarithmen, gewissermaßen verschiedene Systeme geben\begingroup\renewcommand{\thefootnote}{*)}\footnote{Der Unterschied besteht allerdings darin, dass es bei den Logarithmen unendlich viele Systeme gibt, hier dagegen nur so viele wie primitive Wurzeln. Kongruente Basen erzeugen offenbar dasselbe System.}\endgroup.
\SourceLeaf{61}Sehen wir, wie diese zusammenhängen. Seien $a,b$ zwei primitive Wurzeln und $m$ eine weitere Zahl. Wählt man $a$ als Basis, so sei der Index von $b$ $\equiv\beta$ und der Index von $m$ $\equiv\mu\,(\operatorname{mod.}p-1)$. Wählt man dagegen $b$ als Basis, so sei der Index von $a$ $\equiv\alpha$ und der von $m$ $\equiv\nu\,(\operatorname{mod.}p-1)$. Dann ist $\alpha\beta\equiv1\,(\operatorname{mod.}p-1)$. Denn $a^\beta\equiv b$, also nach Voraussetzung $a^{\alpha\beta}\equiv b^\alpha\equiv a\,(\operatorname{mod.}p)$, und damit $\alpha\beta\equiv1\,(\operatorname{mod.}p-1)$. Ebenso findet man $\nu\equiv\alpha\mu$ und $\mu\equiv\beta\nu\,(\operatorname{mod.}p-1)$. Hat man also eine Indextafel für die Basis $a$, so lässt sie sich leicht in eine Tafel zur Basis $b$ umwandeln. Ist nämlich bei der Basis $a$ der Index von $b$ $\equiv\beta$, so ist bei der Basis $b$ der Index von $a$ $\equiv\frac{1}{\beta}\,(\operatorname{mod.}p-1)$. Multipliziert man sämtliche Indizes der Tafel mit dieser Zahl, so erhält man alle Indizes zur Basis $b$.

\Article{70}
Eine gegebene Zahl kann zwar verschiedene Indizes erhalten, wenn man verschiedene primitive Wurzeln als Basis nimmt; doch haben sie alle mit $p-1$ denselben größten gemeinsamen Teiler. Ist nämlich bei der Basis $a$ der Index der gegebenen Zahl gleich $m$ und bei der Basis $b$ gleich $n$, und nimmt man an, die größten gemeinsamen Teiler dieser Indizes mit $p-1$, nämlich $\mu,\nu$, seien ungleich, so ist einer größer, etwa $\mu>\nu$. Damit teilt $\mu$ die Zahl $n$ nicht. Bezeichnet man aber den Index von $a$ bei der Basis $b$ mit $\alpha$, so gilt nach dem vorigen Artikel $n\equiv\alpha m\,(\operatorname{mod.}p-1)$. Also teilt $\mu$ auch $n$, \emph{entgegen der Annahme.}

Dass dieser größte gemeinsame Teiler der Indizes einer gegebenen Zahl und von $p-1$ nicht von der Basis abhängt, ergibt sich auch daraus, dass er gleich $\frac{p-1}{t}$ ist. Dabei bezeichnet $t$ den Exponenten, zu dem die Zahl gehört, um deren Indizes es geht. Ist nämlich der Index für irgendeine Basis gleich $k$, so ist $t$ die kleinste Zahl, mit der man $k$ multiplizieren muss, um ein Vielfaches von $p-1$ zu erhalten, abgesehen von null (siehe Artikel 48, 58). Sie ist also der kleinste Wert von $\frac{0}{k}\,(\operatorname{mod.}p-1)$ außer null. Dass dieser dem größten gemeinsamen Teiler der Zahlen $k$ und $p-1$ gleich ist, ergibt sich ohne Schwierigkeit aus Artikel 29.

\Article{71}
Ferner lässt sich leicht zeigen, dass man die Basis stets so wählen kann, dass eine zum Exponenten $t$ gehörige Zahl jeden vorgegebenen Index erhält, dessen größter gemeinsamer Teiler mit $p-1$ den Wert $=\frac{p-1}{t}$ hat. Bezeichnen wir diesen Teiler kurz mit $d$. Der gewünschte Index sei $\equiv dm$, und wenn eine beliebige
\SourceLeaf{62}primitive Wurzel $a$ als Basis genommen wird, sei der Index der gegebenen Zahl $\equiv dn$. Dann sind $m,n$ zu $\frac{p-1}{d}$ oder $t$ teilerfremd. Ist nun $\varepsilon$ ein Wert von $\frac{dn}{dm}\,(\operatorname{mod.}p-1)$ und zugleich zu $p-1$ teilerfremd, so ist $a^\varepsilon$ eine primitive Wurzel. Nimmt man sie als Basis, so erhält die gegebene Zahl den Index $dm$, wie verlangt (denn $a^{\varepsilon dm}\equiv a^{dn}\equiv$ der gegebenen Zahl). Dass aber $\frac{dn}{dm}\,(\operatorname{mod.}p-1)$ zu $p-1$ teilerfremde Werte hat, beweist man wie folgt. Dieser Ausdruck ist mit $\frac{n}{m}\,(\operatorname{mod.}\frac{p-1}{d})$ oder $\frac{n}{m}\,(\operatorname{mod.}t)$ gleichwertig; siehe Artikel 31, 2. Alle seine Werte sind zu $t$ teilerfremd. Hätte nämlich ein Wert $e$ einen gemeinsamen Teiler mit $t$, so müsste dieser auch $me$ teilen und damit auch $n$, dem $me$ nach $t$ kongruent ist. Das widerspricht der Voraussetzung, dass $n$ zu $t$ teilerfremd ist. Teilen also sämtliche Primteiler von $p-1$ auch $t$, so sind \emph{alle} Werte von $\frac{n}{m}\,(\operatorname{mod.}t)$ zu $p-1$ teilerfremd, und ihre Anzahl ist $=d$. Enthält dagegen $p-1$ noch andere Primteiler $f,g,h$ usw., die $t$ nicht teilen, so wähle man irgendeinen Wert und schreibe $\frac{n}{m}\,(\operatorname{mod.}t)\equiv e$. Da nun sämtliche Zahlen $t,f,g,h$ usw. paarweise teilerfremd sind, lässt sich eine Zahl $\varepsilon$ finden, die nach $t$ der Zahl $e$, nach $f,g,h$ usw. dagegen beliebigen zu diesen jeweils teilerfremden Zahlen kongruent ist (Artikel 32). Eine solche Zahl ist durch keinen Primfaktor von $p-1$ teilbar und daher zu $p-1$ teilerfremd, wie verlangt. Schließlich ergibt sich aus der Kombinationslehre ohne Schwierigkeit die Anzahl $=\frac{p-1\,.\,f-1\,.\,g-1\,.\,h-1.\ \mathrm{usw.}}{t\,.\,f\,.\,g\,.\,h\,.\,\mathrm{usw.}}$ dieser Werte. Damit diese Abschweifung nicht zu umfangreich wird, lassen wir den Beweis weg, da er für unseren Zweck nicht besonders nötig ist.

\begin{center}\emph{Basen für besondere Zwecke.}\end{center}

\Article{72}
Welche primitive Wurzel als Basis gewählt wird, ist im Allgemeinen völlig beliebig. Manchmal bieten aber bestimmte Basen besondere Vorteile gegenüber anderen. In Tafel I haben wir stets die Zahl $10$ als Basis gewählt, wenn sie eine primitive Wurzel war. Andernfalls haben wir die Basis immer so bestimmt, dass der Index von $10$ möglichst klein wurde, \emph{das heißt} $=\frac{p-1}{t}$. Dabei bezeichnet $t$ den Exponenten, zu dem $10$ gehört. Was wir dadurch gewinnen, zeigen wir in Abschnitt VI, wo dieselbe Tafel noch für andere Zwecke verwendet wird. Da aber, wie aus dem vorigen Artikel hervorgeht, auch hier noch eine Wahl bleiben kann, haben wir eine feste Regel getroffen: Von allen primitiven Wurzeln, die das Verlangte leisten, haben wir stets die \emph{kleinste} als Basis gewählt. So hat für $p=73$, wobei $t=8$ und $d=9$ ist, $a^\varepsilon$ genau $\frac{72\,.\,2}{8\,.\,3}$, \emph{das heißt} $6$ Werte, nämlich $5,14,20,28,39,40$. Wir haben also den kleinsten, $5$, als Basis genommen.
